\chapter{Introduction}\label{ch:intro}

% ------------------------------------------------------------
% Opening: why economics for AI alignment.
% Lifted (adapted) from marketic_agents_project_intro.org and
% shortform notes; glue composed for the thesis is draft-marked.
% ------------------------------------------------------------

\begin{draftblock}
Economics is the study of aligning selfish agents to a social goal; AI alignment is the problem of ensuring that artificial agents pursue their principals' goals. This thesis takes the position that this is not a loose analogy but a productive research programme: that the conceptual toolkit of economics --- prices, markets, information asymmetry, mechanism design, belief elicitation --- is the right language for several of the central problems of AI alignment, and conversely that AI agents (in particular Large Language Models) resolve long-standing obstructions to economic mechanisms that were previously impossible with human participants alone.
\end{draftblock}

% Lifted from marketic_agents_project_intro.org (opening).
Markets are an example of real-life \emph{Aligned Superintelligence}, and they serve as a source for intuition-pumping for alignment. Markets serve as an example of ``already-existing (aligned) AGI'', so they serve as a nice setting for us to play with \emph{concepts of intelligence and alignment} --- even if we don't want to build market-based agents, markets serve as a valuable intuition pump for the AI agents we do build.

\section{Markets and intelligence}\label{intro:sec:markets-intelligence}

% Lifted from the year-2 PhD report ("Markets and intelligence"), Introduction.
Before neural networks won the mandate of heaven, an AI research paradigm that had shown considerable promise was that of \emph{market-based architectures}, i.e. AI agents that determine their output based on some internal market-based mechanism \citep{kweeMarketBasedReinforcementLearning2001, changDecentralizedReinforcementLearning2020}.

More broadly speaking, there are general intuitive arguments to imagine a close analogy between markets and AI. Philosophers and psychologists have long pondered multi-agent models of the mind \citep{minskySocietyMind1988}; moreover, one may imagine that ``any'' machine learning task could in principle be solved by a market of agents solving sub-tasks with their individual reward set by the sale value of their output. There is work suggesting that markets can capture some notion of \emph{bounded rationality}, e.g. the \emph{Boundedly Rational Inductive Agent} (BRIA) \citep{oesterheldTheoryBoundedInductive2023} and \emph{Algorithmic Bayesian Epistemology} \citep{neymanAlgorithmicBayesianEpistemology2024} (intuitively, the efficient market hypothesis may be imagined as expressing that markets are efficient conditional on the \emph{algorithmic information} available). In some sense, markets ``aggregate'' the intelligence or capacities of their individual participants.\footnote{See e.g. Hayek on the role of markets in aggregating information \citep{hayekEconomicsKnowledge1937} -- or the famous parable ``I, Pencil'' \citep{leonardereadPencilMyFamily1958}: ``... no one person, no matter how smart, could create from scratch a small, everyday pencil [yet the market makes over a billion of them each year] ...''. There is also some empirical work on emergent intelligent behaviour in markets comprised of zero-intelligence traders \citep{godeAllocativeEfficiencyMarkets1993, schwartzHowMuchIrrationality2008, jamalSimpleAgentsIntelligent2015}}

Research at the intersection of markets and AI may be split into two overarching categories\footnote{There is substantial overlap between these categories as one may imagine a market ensemble of AI agents as itself a ``super-AI'', which has by construction a market structure}:

\begin{itemize}
\item \textbf{Markets \emph{as} AI,} i.e. AI models and agents that work based on internal market-based algorithms, i.e. with intelligence arising as an ``emergent property'' of several simple or bounded agents interacting economically. This includes market-based RL (see \cref{ch:market-rl}), market architectures for logical uncertainty \citep{garrabrantLogicalInduction2020}, market architectures for bounded rationality \citep{oesterheldTheoryBoundedInductive2023} and mechanism design for reward \citep{conitzerPositionSocialChoice2024, geAxiomsAIAlignment2024a}.
\item \textbf{Markets \emph{of} AI,} i.e. markets comprised of AI participants. This includes economic simulations with AI populations e.g. \citep{zhengAIEconomistOptimal2021, liEconAgentLargeLanguage2024a}; AI forecasters participating on prediction markets (see \cref{ch:consistency}) and decentralized ``intelligence markets'' such as BitTensor \citep{raoBitTensorPeertoPeerIntelligence2021}.
\end{itemize}

% Lifted from marketic_agents_project_intro.org ("Analogies between markets and agents").
Markets and agents do the ``same kinds of things'' (i.e. have the same ``type signature''): they compute \emph{decisions} (in the form of goods allocation) and \emph{beliefs} (via e.g. prediction markets). Both markets and \emph{boundedly rational} agents are only ``rational given the available algorithmic information''. Markets have \emph{self-awareness}: they can contain prediction markets for the market's own properties, and can also self-improve, because the algorithmic functionality of markets is itself provided by goods (e.g. ``the labour supply of arbitrageurs'') whose supply can be adjusted and optimized. Abstraction is like \emph{economies of scale}; \emph{transaction costs} are like \emph{computational costs}. Markets are \emph{general intelligence}: they can be asked (paid) to accomplish \emph{any} task --- and they are intelligent about managing this generality, forming abstractions (factor goods) that will be useful for the distribution of tasks they expect to be asked to accomplish.

\section{The alignment problem, economically}\label{intro:sec:alignment}

\begin{draftblock}
If markets are the economist's model of distributed intelligence, the alignment problem is the economist's principal--agent problem at scale. It is useful to decompose the standard rebuttal to ``why can't we just set the right reward function?'' into three levels:
\end{draftblock}

% Lifted from alignment_summary.org ("Summary").
\begin{itemize}
\item \textbf{The problem of specification} --- How do you communicate what you want?
\item \textbf{The problem of information asymmetry (outer alignment)} --- Giving reward signals to a superhuman AI for its behaviour is fundamentally subject to ``information asymmetry'': you are not good enough to correctly evaluate the AI's behaviour. How do you ensure your reward signals reflect your \emph{extrapolated volition} (i.e. incorporating all the information in the world) rather than your immediate, superficial volition?
\item \textbf{The problem of being weak and pathetic (inner alignment)} --- To a superhuman AI, your reward function is a \emph{joke}. You have a being of unfathomable intelligence and you are trying to edit its brain bit-by-bit with something called ``backpropagation''. The reward function is not the utility function; it is your puny and laughable attempt to shape a superhuman AI into something desirable.
\end{itemize}

% Lifted from ways_to_think_alignment.org (A).
\subsection{Buying under information asymmetry}\label{intro:sec:infoasym}

The basic problem with RLHF is that it fundamentally relies on a human's ability to judge the correctness or value of a (potentially superhuman) AI's outputs. In other words, the AI is trained on the human supervisor's immediate, superficial volition, rather than on her extrapolated volition.

The standard measurement device for human values is the ``market'' --- you put something on the market, and see how much it sells for. But markets, too, only price things correctly when assuming \emph{perfect information} --- if you are unable to know (or verify/trust, if told) something about a product, e.g. ``how many years will this battery last?'' ``what is the polyphenolic content of this olive oil?'' ``is this item adulterated?'', it doesn't affect your decision to buy, so the producer has no reason to optimize on it, so he makes the cheapest possible product; you learn to expect this and revise down your valuation of it. More generally, fixing this is the goal of any \emph{scalable oversight} mechanism.

% Lifted from ways_to_think_alignment.org (B), abridged.
\subsection{Utility $\neq$ reward}\label{intro:sec:utility-reward}

As a simple demonstrative example: the \emph{reward function} that trains (trained) humans is ``survival and reproduction in the (current) (ancestral) environment''. But this is not the \emph{utility function}, i.e. what we care about in life. Your specified reward function is not the \emph{true} reward function given by nature: ``supplying a reward function'' is just your attempt at modifying the natural environment --- specifically by constantly editing the AI by tweaking its parameters based on gradient descent. This training process will work \emph{only as long as that is the dominant selective pressure on the AI}.

The market analogy is instructive: the market selects for agents which serve the consumers well --- such agents gain influence in the market, and as a result the aggregate of agents learns to optimize for social value. But this does not mean you have selected for agents that ``care'' about what you want, just ones which are very good at giving you what you want. The market mechanism does not, in itself, prevent a corporation from simply going rogue; it doesn't even make agents care about money.

\begin{draftblock}
Understanding what reward \emph{is} --- what role it actually plays in producing agency, and what the rational theory of it looks like --- is thus itself a foundational question for alignment, prior to any particular alignment scheme. \Cref{part:agency} of this thesis develops an answer in economic terms: reward as price, and price recursion as the rational theory of reward.
\end{draftblock}

% Lifted from ways_to_think_alignment.org (C), abridged.
\subsection{Prover--verifier games and subjective questions}\label{intro:sec:pv}

There are, one might say, three sorts of information in the world:
\begin{itemize}
\item ``Ideas'': math/science theories and models, inventions, business ideas, solutions to open-ended problems;
\item ``Answers'': math theorems, experimental observations, results of computations;
\item ``Proofs'': math proofs, arguments, evidence, digital signatures, certifications, reputations, signalling.
\end{itemize}
\ldots{} and a reasonable analogy: Ideas~:~search, Answers~:~inference, Proofs~:~\emph{alignment}. Suppose you have a mechanism that bridges the information asymmetry gap between the human and the AI. You still need to ensure that the human \emph{learns} to correctly update on this information --- the AI (the prover) and the human (the verifier) need to learn in tandem, so the former can produce ``proofs'' understandable by the latter. This is never a problem if the human were a \emph{Bayesian} agent --- the AI could simply provide all the relevant evidence to convince the human of the truth. The ability to ``rationalize'' a belief update mechanism as ``just Bayesian lacking real information'' is what characterizes \emph{bounded rationality}.

\begin{draftblock}
A special difficulty arises when the questions on which beliefs must be elicited are not even \emph{verifiable}: when no ground truth will ever arrive to settle the bet. Human values, latent variables, and ``subjective'' judgements are of exactly this character. The economist's elicitation devices --- proper scoring rules and prediction markets --- presuppose eventual resolution; extending them beyond that boundary is the subject of \cref{part:subjective}.
\end{draftblock}

\section{Three levels: overview of the thesis}\label{intro:sec:overview}

\begin{draftblock}
The thesis applies economics to AI alignment at three levels, in decreasing order of ambition and increasing order of concreteness:

\begin{enumerate}
\item \textbf{An economic theory of agency} (\cref{part:agency}). At the deepest level, we take the market not merely as a metaphor for intelligence but as a candidate \emph{theory} of it. \Cref{ch:market-rl} studies market-based architectures in reinforcement learning and beyond: agents whose internal credit assignment is a market, with the state factored into ``goods'' and subtasks priced by internal auction. \Cref{ch:agency} develops the accompanying theory --- that \emph{price recursion is the rational theory of reward}, with reward understood as hyperstitional information propagated backwards through an economy of computations --- together with the \texttt{marketgrad} library implementing it.
\item \textbf{Information asymmetry and scalable oversight} (\cref{part:oversight}). At the level of the alignment problem proper, we identify the outer alignment/scalable oversight problem with the economist's problem of information asymmetry between principal and agent. \Cref{ch:infonomy} proposes and analyzes \emph{recursive information markets} as a mechanism for extrapolating the volition of a human rater; \cref{ch:solib} constructs a general benchmark for evaluating scalable oversight protocols.
\item \textbf{Eliciting beliefs on subjective questions} (\cref{part:subjective}). Any of the above requires eliciting honest beliefs --- and the questions that matter for alignment are often ones on which no ground truth ever arrives. \Cref{ch:nonvf} constructs scoring mechanisms for claims that are neither verifiable nor falsifiable; \cref{ch:consistency} develops consistency checks --- in particular arbitrage-based metrics --- for evaluating and improving language-model forecasters even before resolution.
\end{enumerate}
\end{draftblock}

% Adapted from the year-2 PhD report: the fundamental research questions.
Over the course of this research, the following fundamental questions kept recurring, and organize the chapters that follow:

\begin{enumerate}
\item \emph{Can we design a reinforcement learning agent based on an internal market-based algorithm?} Would doing so have any advantages?
\item Some works have noted a suggestive market analogy in neural networks \citep{davisMappingClassifierSystems1988, schmidhuberLocalLearningAlgorithm1989} and in backpropagation \citep{wentworthCompetitiveMarketsDistributed2018}. \emph{Can we make this more precise, and use it to construct a market-based algorithm for supervised learning applications?}
\item Markets may be thought of as measuring devices for people's values and beliefs -- \emph{can we develop some ``AIs as producers, humans as consumers'' framework for giving human feedback to AIs, i.e. replacing naive RLHF?}
\item Prediction markets can easily give probabilities for verifiable or falsifiable (VF) events. But many questions we are interested in are not VF -- they exist only in the ``latent space''. \emph{Can we design scoring rules to incentivize usefully reporting probabilities on latent space variables?} The resulting economy may be seen as an AI that has learned some internal beliefs, and may be valuable for interpretability research \citep{christianoElicitingLatentKnowledge2021}.
\end{enumerate}

\begin{draftblock}
Questions 1 and 2 are the subject of \cref{part:agency}; question 3 of \cref{part:oversight}; question 4 of \cref{part:subjective}.
\end{draftblock}
